\chapter{Building Servers}
\label{ch:building-servers}

\epigraph{``Do, or do not. There is no try.''}{Yoda, \emph{The Empire Strikes
Back} (1980)}

Yoda is talking about commitment, and he is also, accidentally, describing
error handling. A tool either did the thing or it did not, and the worst
possible outcome is a response that leaves the model unsure which. Half the
craft of building a good server is making sure every response says
unambiguously which of those two happened.

This chapter builds one from an empty directory. Not a toy. The risk server that
ships with this book, with everything real: MRTR approvals, TTL'd resources,
Tasks, a dual-era bridge, and contract tests that pass.

\section{Anatomy of a 2026 server}

Start with what is \emph{not} here, because it is more than you expect.

\begin{itemize}[leftmargin=1.6em, itemsep=2pt]
\item No \texttt{initialize} handler.
\item No session object.
\item No connection state.
\item No \texttt{ping}.
\item No \texttt{logging/setLevel}.
\item No \texttt{resources/subscribe}.
\end{itemize}

What remains is a router:

\begin{py}
class Server:
    """A stateless MCP server.

    Note what this class does not have. No `initialize`. No session object.
    No per-connection state of any kind. A `Server` is a pure function from
    request to response, which is why you can run twenty of them behind a
    round-robin load balancer and never think about stickiness again.
    """
\end{py}

The whole dispatch path fits on a page, and it is worth reading closely because
four of its decisions come straight from the specification:

\begin{py}
def handle(self, message, *, auth=None, emit_progress=None):
    try:
        jsonrpc.validate_request(message)
    except errors.McpError as exc:
        return jsonrpc.error_response(message.get("id"), exc)

    method = message["method"]
    params = message.get("params") or {}
    req_id = message.get("id")

    if jsonrpc.is_notification(message):
        self.handle_notification(method, params)
        return None                                    # (1)

    try:
        ctx = parse_request_context(                   # (2)
            method, params, request_id=req_id, auth=auth,
            supported_versions=self.supported_versions)
        ctx.emit_progress = emit_progress
        handler = self._methods.get(method)
        if handler is None:
            raise errors.MethodNotFound(method)

        result = handler(ctx)
        result.setdefault("resultType", jsonrpc.RESULT_COMPLETE)   # (3)
        meta = result.setdefault("_meta", {})
        meta.setdefault(KEY_SERVER_INFO, self.info.to_json())      # (4)
        return jsonrpc.result_response(req_id, result)

    except errors.McpError as exc:
        return jsonrpc.error_response(req_id, exc)
    except Exception as exc:  # a handler bug must not take the process down
        return jsonrpc.error_response(
            req_id, errors.InternalError(f"{type(exc).__name__}: {exc}"))
\end{py}

\begin{enumerate}[leftmargin=1.6em, itemsep=3pt]
\item Notifications get no response. Returning one is a protocol violation and
some clients will treat the unexpected message as a fatal error.
\item Every request is validated against \texttt{\_meta} before anything else
happens. No handshake means no other opportunity.
\item \texttt{resultType} defaults to \texttt{complete} so handlers never have to
think about it.
\item The server identifies itself in every result, as the specification
recommends.
\end{enumerate}

And the bare \texttt{except} at the bottom is deliberate. A handler that raises
\texttt{KeyError} should produce a \texttt{-32603} for that one request, not take
down a process serving forty others.

\section{Registering the surface}

Meridian uses decorators, because they keep the tool next to its schema and its
description, which is where they belong.

\begin{py}
@server.tool(
    "assess_account_risk",
    "Score one commercial account for credit risk. Returns a 1-99 score, a "
    "band, and the weighted factors behind it.",
    {
        "type": "object",
        "properties": {
            "accountId": {
                "type": "string",
                "description": "Account identifier, e.g. ACC-1042",
                "pattern": "^ACC-[0-9]{4}$",
            },
            "exposureUsd": {
                "type": "number", "minimum": 0,
                "description": "Proposed exposure. Above 5,000,000 an "
                               "approver is required before scoring.",
            },
            "region": {
                "type": "string",
                "enum": ["us-east", "us-west", "eu-central", "eu-west",
                         "apac-south"],
                "description": "Routing hint. Mirrored into an HTTP header "
                               "so gateways can route without reading the body.",
                "x-mcp-header": "Region",
            },
        },
        "required": ["accountId"],
    },
    output_schema={ ... },
    annotations={"readOnlyHint": True, "idempotentHint": True},
    ui_resource_uri="ui://meridian/risk-dashboard",
)
def assess_account_risk(ctx: RequestContext):
    ...
\end{py}

Four things in that declaration earn their place, and each connects to an
earlier chapter.

The \texttt{pattern} catches a malformed id at the boundary, so the handler never
sees it. The \texttt{description} on \texttt{exposureUsd} tells the model about
the approval threshold, so it can warn the user \emph{before} triggering an MRTR
round trip that costs 375 milliseconds. The \texttt{x-mcp-header} lets a gateway
route by region without parsing the body. And the annotations let a host
auto-approve, because scoring an account changes nothing.

\subsection{Validate at the edge, then trust}

Meridian validates arguments against the input schema before calling the
handler:

\begin{py}
def _tools_call(self, ctx: RequestContext) -> dict:
    name = ctx.params.get("name")
    tool = self._tools.get(name)
    if tool is None:
        raise errors.InvalidParams(f"Unknown tool: {name}")

    validate_against_schema(ctx.arguments, tool.input_schema,
                            where=f"{name} arguments")
    result = tool.handler(ctx)
    return normalise_tool_result(result, tool)
\end{py}

Which means handlers can write \texttt{ctx.arguments["accountId"]} without a
guard, because a request that got this far has one. That is a small quality-of-
life thing that compounds across forty tools.

\subsection{Let handlers return whatever is natural}

A handler returning a string, a dict, or a full envelope should all work. That is
one function:

\begin{py}
def normalise_tool_result(result, tool) -> dict:
    """Let handlers return a string, a dict, or a full result envelope."""
    if isinstance(result, str):
        return text_result(result)
    # Already an InputRequiredResult or a full envelope: pass it through.
    if "resultType" in result or "content" in result:
        return result
    # A bare dict is structured content. Mirror it into text for the model,
    # per the backwards-compatibility guidance in the spec.
    return text_result(json.dumps(result, separators=(",", ":")),
                       structured=result)
\end{py}

That last branch implements a specification recommendation quietly: a tool
returning structured content should also return the serialised JSON as text, so
clients that do not read \texttt{structuredContent} still work.

\section{State without sessions}

Chapter~\ref{ch:architecture} introduced the handle pattern: a server that needs
to remember something between calls mints an identifier, returns it, and requires
it back. It also gave the four rules that keep one safe. This section is about
the part that chapter left out, which is what the code looks like.

The whole mechanism on the server side is a dictionary and an authorization
check:

\begin{py}
@server.tool("add_item", "Add an item to an open basket.", ADD_ITEM_SCHEMA)
def add_item(ctx: RequestContext):
    basket_id = ctx.arguments["basket_id"]
    basket = BASKETS.get(basket_id)

    # A handle names an object. It does not grant access to one, so the
    # caller's authorization is re-checked here on every single call.
    if basket is None or not can_access(ctx.auth, basket):
        return tool_error(
            f"Basket {basket_id} does not exist or has expired. "
            "Create a new one with create_basket.")

    basket.items.append(ctx.arguments["sku"])
    basket.touched_at = time.time()
    return {"structuredContent": {"basket_id": basket_id,
                                  "itemCount": len(basket.items)}}
\end{py}

Three details in that handler are doing real work.

The authorization check and the existence check produce the \emph{same} error.
Distinguishing ``this basket is not yours'' from ``this basket does not exist''
tells an attacker which identifiers are real, which is how you turn a handle
space into an enumeration oracle.

The error is a tool execution error, so the model reads it. And it names the
recovery: create a new one. Chapter~\ref{ch:tools} argued for this generally;
here it is the difference between a model that recovers from an expired basket
and a model that retries a dead identifier until the step budget stops it.

And \texttt{touched\_at} exists because handles need a retention policy, which
needs to be visible to the model:

\begin{py}
@server.tool(
    "create_basket",
    "Create a shopping basket and return its id. Baskets expire after 24 "
    "hours of inactivity; pass the id to add_item and checkout.",
    {"type": "object", "additionalProperties": False},
)
\end{py}

That description is the only place the model can learn the lifetime. It is not
in the schema, it is not in the protocol, and there is nowhere else to put it.

\section{Idempotency, because retries are routine now}

Chapter~\ref{ch:transport} removed stream resumability. A broken response stream
means the client re-issues with a new id, and it has no way to know whether the
first attempt reached you.

So retries are normal operation. Which means any tool with side effects needs an
idempotency key, and the protocol will not provide one. It is an ordinary
argument, which is exactly why it works everywhere with no negotiation:

\begin{py}
IDEMPOTENCY_KEY_SCHEMA = {
    "type": "string",
    "minLength": 8,
    "maxLength": 128,
    "description": (
        "Client-generated unique id for this intent. Repeating a key within "
        "24 hours returns the original result rather than performing the "
        "action again. Generate a fresh key per intent; reuse it on retries."
    ),
}
\end{py}

That description is load-bearing. The model is the thing generating the key, and
the two sentences tell it the rule that makes the mechanism work: one key per
intent, the same key on every retry of that intent. A model that generates a
fresh key per attempt has an argument that does nothing.

The naive server side is a dictionary lookup, and the naive version is wrong in
four ways. Meridian's is worth reading for the four things it does past the
lookup.

\textbf{The key is not the cache key.} Two different callers may generate the
same key, so the record is scoped to the principal. And one caller may reuse a
key with different arguments, which is a client bug: returning the first
operation's result would quietly answer a question nobody asked.

\begin{py}
# A replay. Reusing one key for two different requests is a client bug
# worth surfacing loudly, because the alternative is returning one
# operation's result as though it were another's.
if record.fingerprint != digest:
    raise errors.InvalidParams(
        "idempotencyKey was already used with different arguments")
\end{py}

\textbf{A concurrent duplicate must wait, not proceed.} Two retries can be in
flight at once, and a check-then-act on a plain dictionary lets both find it
empty. The record is inserted before the work starts, and the second arrival
blocks on the first:

\begin{py}
if not record.done.wait(wait_timeout):
    raise errors.InternalError(
        "a request with this idempotencyKey is still in flight")
\end{py}

\textbf{A failure must not be remembered.} Storing an exception as though it
were a result makes a transient downstream outage permanent for the life of the
key, which is 24 hours of a client that can never succeed:

\begin{py}
except BaseException:
    # A failure must not be remembered as a result, or the caller
    # can never retry it. Drop the record and let the next attempt
    # start clean.
    with self._lock:
        self._records.pop(slot, None)
    record.done.set()
    raise
\end{py}

\textbf{Records expire.} The retention window is the promise you are making
about how long a retry works. Sweep afterwards, or the store grows for the life
of the process.

The same caveat applies here as to the task store in Chapter~\ref{ch:tasks}:
in-memory works in a book and fails behind a load balancer, where the retry
lands on a replica that has never heard of the key and cheerfully does the work
again.

\section{Configuration}

Environment variables for secrets. Nothing exotic, but two details that bite.

\begin{py}
SEALER = StateSealer(
    secret=os.environ["MERIDIAN_STATE_SECRET"].encode(),
    ttl_seconds=900,
)
\end{py}

\textbf{The signing secret must be fleet-wide.} Meridian's default generates a
random one per process, which is correct for a single-process demo and
catastrophic behind a load balancer: a retry landing on replica B cannot open
state sealed by replica A, and the user sees an approval that mysteriously fails
about half the time.

\textbf{Declare extensions honestly.} Advertising \texttt{tasks} without
implementing \texttt{tasks/get} produces a client that polls a method that
returns \texttt{-32601} forever.

\begin{py}
def declare_extension(self, ident: str, settings: dict | None = None) -> None:
    """Advertise an extension. Declaring one you do not implement is lying,
    and the client will believe you."""
\end{py}

\section{Capabilities should be derived, not declared}

A small design decision that removes a whole class of bug. Meridian computes its
capabilities from what is actually registered:

\begin{py}
def capabilities(self) -> ServerCapabilities:
    caps = ServerCapabilities(extensions=dict(self._extensions))
    if self._tools:
        caps.tools = {"listChanged": True} if self.list_changed else {}
    if self._resources or self._templates:
        res = {}
        if self.list_changed:
            res["listChanged"] = True
        if self.subscribe:
            res["subscribe"] = True
        caps.resources = res
    if self._prompts:
        caps.prompts = {"listChanged": True} if self.list_changed else {}
    return caps
\end{py}

A hand-maintained capability list drifts from reality the first time somebody
deletes the last prompt and forgets the declaration. A derived one cannot.

\section{Filtering by authorization, not by connection}

The specification draws a line here that is easy to misread. A server's tool list
\textbf{must not} vary per connection. It \textbf{may} vary by the authorization
presented on the request.

Those sound similar and are opposites. Credentials arrive with each request, so
filtering on them is stateless. Connection identity is exactly the thing this
revision deleted.

\begin{py}
def visible_tools(self, ctx: RequestContext) -> list[Tool]:
    """Override to filter by scope. The set may vary by authorization,
    never by connection."""
    return sorted(self._tools.values(), key=lambda t: t.name)
\end{py}

A scope-aware subclass:

\begin{py}
class ScopedRiskServer(Server):
    """The risk server, with per-tool scopes and row-level account access."""

    def visible_tools(self, ctx: RequestContext) -> list[Tool]:
        """Override to filter by scope. The set may vary by authorization,
        never by connection."""
        scopes = set((ctx.auth or {}).get("scopes", []))
        return [t for t in super().visible_tools(ctx)
                if not REQUIRED_SCOPE.get(t.name)
                or REQUIRED_SCOPE[t.name] in scopes]
\end{py}

Note that \texttt{super()} sorts first, so the filtered list is still
deterministic, which keeps the prompt-cache property from
Chapter~\ref{ch:tools}. Filtering after sorting preserves order; sorting after
filtering would too, but doing it in the wrong place is how people lose it.

\section{Wiring the transports}

The server is transport-agnostic. The transports are thin.

\begin{py}
def main(argv=None) -> int:
    ap = argparse.ArgumentParser(description="Meridian risk server")
    ap.add_argument("--http", type=int, metavar="PORT")
    ap.add_argument("--fat", action="store_true")
    args = ap.parse_args(argv)

    server = build_server(fat_catalogue=args.fat)
    tasks_ext.install(server)

    if args.http:
        http = StreamableHttpServer(server, port=args.http)
        print(f"meridian-risk on {http.url}", file=sys.stderr, flush=True)
        http.start()
        ...
    StdioServerTransport(server).serve_forever()
\end{py}

That \texttt{file=sys.stderr} is not decoration. On stdio, \texttt{stdout} carries
MCP messages and nothing else, and a startup banner on \texttt{stdout} corrupts
the stream before the first request. It is the single most common way a new
stdio server fails, and the error the client reports gives no hint of the cause.

\subsection{Serving both eras}

Chapter~\ref{ch:architecture} explained why you need this in 2026. The wiring is
one wrapper:

\begin{py}
server = build(args.server, **kwargs)
endpoint = server if args.modern_only else LegacyBridge(server)

if args.http:
    http = StreamableHttpServer(endpoint, port=args.http)
    ...
StdioServerTransport(endpoint).serve_forever()
\end{py}

Dual-era is the default because that is what actually connects to shipping
clients. \texttt{-{}-modern-only} exists so you can watch a strict server refuse
a handshake, which is a useful thing to see once.

\section{Contract tests}

A server without contract tests is a server that conforms to the specification
today.

Meridian has 191 tests, and the ones in this section are the ones I would write
first for any new server. They assert on the wire, not on convenience objects,
because the wire is what another implementation has to interoperate with.

\subsection{The envelope}

\begin{py}
def test_request_without_protocol_version_is_rejected(self):
    resp = self.server.handle({
        "jsonrpc": "2.0", "id": 1, "method": "tools/list",
        "params": {"_meta": {KEY_CLIENT_CAPABILITIES: {}}},
    })
    self.assertEqual(resp["error"]["code"], errors.INVALID_PARAMS)
    self.assertIn("protocolVersion", resp["error"]["message"])

def test_unsupported_version_lists_what_is_supported(self):
    resp = self.server.handle(request("tools/list", version="1999-01-01"))
    self.assertEqual(resp["error"]["code"], errors.UNSUPPORTED_PROTOCOL_VERSION)
    self.assertIn(PROTOCOL_VERSION, resp["error"]["data"]["supported"])
\end{py}

That second one matters more than it looks. The error must carry a
\texttt{supported} list, because that list is how a client renegotiates without
a handshake. An error that just says ``unsupported'' leaves the client with
nothing to do.

\subsection{Statelessness itself}

\begin{py}
def test_two_requests_share_no_state(self):
    """The second request must not benefit from the first having happened."""
    first = self.server.handle(request("tools/list", req_id=1))
    second = self.server.handle(request("tools/list", req_id=2))
    self.assertEqual(first["result"]["tools"], second["result"]["tools"])
    # And a request with no _meta still fails, even after a good one.
    third = self.server.handle(
        {"jsonrpc": "2.0", "id": 3, "method": "tools/list", "params": {}})
    self.assertIn("error", third)
\end{py}

The second half is the real test. It catches a server that ``helpfully''
remembers the last known-good capabilities, which is a plausible optimisation
and a specification violation.

\subsection{Caching hints on all six operations}

\begin{py}
def test_all_six_cacheable_methods_carry_hints(self):
    for method, params in [
        ("server/discover", {}), ("tools/list", {}), ("prompts/list", {}),
        ("resources/list", {}), ("resources/templates/list", {}),
        ("resources/read", {"uri": "test://doc"}),
    ]:
        with self.subTest(method=method):
            result = server.handle(request(method, params))["result"]
            self.assertIn("ttlMs", result, f"{method} is missing ttlMs")
            self.assertIn("cacheScope", result)
            self.assertGreaterEqual(result["ttlMs"], 0)
            self.assertIn(result["cacheScope"], ("public", "private"))
\end{py}

\texttt{subTest} means one failing method reports as one failure, so you see all
six instead of stopping at the first.

\subsection{Pagination that actually terminates}

\begin{py}
def test_pages_cover_everything_exactly_once(self):
    server = tiny_server(page_size=7)
    for i in range(30):
        server.add_tool(_noop_tool(f"t{i:02d}"))

    seen, cursor, pages = [], None, 0
    while True:
        params = {"cursor": cursor} if cursor else {}
        result = server.handle(request("tools/list", params))["result"]
        seen.extend(t["name"] for t in result["tools"])
        pages += 1
        cursor = result.get("nextCursor")
        if not cursor:
            break
        self.assertLess(pages, 20, "pagination did not terminate")

    self.assertEqual(len(seen), 31)              # 30 plus the built-in `add`
    self.assertEqual(len(seen), len(set(seen)))  # no duplicates
\end{py}

Three properties in one test: it terminates, it covers everything, and it
duplicates nothing. Off-by-one cursor bugs fail one of the three and are
otherwise invisible.

\subsection{The error taxonomy}

\begin{py}
def test_unknown_tool_is_a_protocol_error(self):
    resp = tiny_server().handle(request("tools/call", {"name": "nope"}))
    self.assertIn("error", resp)

def test_bad_business_input_is_a_tool_execution_error(self):
    resp = build_server().handle(request(
        "tools/call",
        {"name": "assess_account_risk", "arguments": {"accountId": "ACC-9999"}}))
    self.assertNotIn("error", resp)
    self.assertTrue(resp["result"]["isError"])
    # And it must be actionable, not just "invalid".
    self.assertIn("ACC-1000", resp["result"]["content"][0]["text"])
\end{py}

That last assertion is unusual and I recommend it. It tests the \emph{quality} of
an error message, which is normally the kind of thing nobody tests and everybody
regrets. If a refactor turns your helpful message into ``not found'', the test
fails.

\section{The build, start to finish}

Meridian's risk server, in order:

\begin{enumerate}[leftmargin=1.6em, itemsep=3pt]
\item \textbf{Construct}, with instructions written for a model to read.
\item \textbf{Three tools}, task-shaped: one account, many accounts, and the
long-running one.
\item \textbf{MRTR approval} above the exposure threshold, with sealed state.
\item \textbf{Resources}: templated account summaries (private, 15 min), public
filings (public, 24 h), a model card (public, 7 d), and the App template.
\item \textbf{A versioned prompt.}
\item \textbf{Tasks}, installed as an extension.
\item \textbf{Transports}, dual-era by default.
\end{enumerate}

The whole file is around 500 lines including the HTML template, and the
\texttt{instructions} string is worth calling out because it is the one piece of
prose the model reads before deciding anything:

\begin{py}
instructions=(
    "Credit risk scoring for commercial accounts. Start with "
    "`assess_account_risk`; it returns a score, a band, and the factors "
    "that drove it. For a portfolio view use `rank_portfolio_risk` "
    "rather than calling the per-account tool in a loop."
),
\end{py}

Two sentences, and the second one prevents the most expensive mistake a model
can make against this server. That is a good return on 30 tokens.

\section{Common mistakes}

\begin{table}[htbp]
\centering
\small
\begin{tabularx}{\textwidth}{@{}XX@{}}
\toprule
\thead{Mistake} & \thead{Symptom} \\
\midrule
Printing to \texttt{stdout} on stdio & Client reports a parse error it cannot attribute \\
Per-process signing secret & MRTR retries fail intermittently behind a load balancer \\
Non-deterministic \texttt{tools/list} order & Prompt-cache hit rate silently collapses \\
Hand-maintained capability list & Advertises prompts you deleted last quarter \\
Protocol error for bad business input & Model gives up instead of retrying with a valid id \\
Advertising an unimplemented extension & Client polls a \texttt{-32601} forever \\
In-memory task store behind a balancer & ``Unknown task'' on roughly $1 - 1/N$ of polls \\
No idempotency key on a mutating tool & A retried stream charges the card twice \\
\bottomrule
\end{tabularx}
\caption{Everything here has happened. Most of them fail silently or
intermittently, which is why they are worth a checklist.}
\label{tab:server-mistakes}
\end{table}

\section{What to remember}

\begin{itemize}[leftmargin=1.6em, itemsep=3pt]
\item A server is a pure function from request to response. No handshake, no
session, no connection state.
\item Validate \texttt{\_meta} first, validate arguments at the edge, then let
handlers trust their inputs.
\item Identify yourself in every result. Default \texttt{resultType} centrally.
\item Cross-call state is an opaque handle passed as an ordinary argument, and
authorized on every use.
\item Retries are routine. Mutating tools need an idempotency key.
\item Derive capabilities from what is registered. Declare extensions only when
you implement them.
\item Filter by authorization, never by connection, and sort before you filter.
\item On stdio, \texttt{stdout} is for MCP messages only.
\item Sign MRTR state with a fleet-wide secret.
\item Write contract tests that assert on the wire, including one that checks
your error messages are actually helpful.
\end{itemize}

Next: the other side of the connection, where the trust boundary lives.
